\section{Introduction}

Time-series transcriptomics answers \emph{when} a response happens, but not \emph{who}
drives it. The Dynamic Regulatory Events Miner (DREM) closes that gap by fitting an
input--output hidden Markov model in which co-expressed genes travel together until a
\emph{bifurcation}, and by attributing each bifurcation to the transcription factors
(TFs) whose targets are enriched along one branch \citep{ernst2007drem}. DREM~2.0 added
model-selection and scoring improvements \citep{schulz2012drem2}, and iDREM extended the
framework to multi-omic inputs and interactive exploration \citep{ding2018idrem}.

The attribution step depends entirely on a second input: a static TF--gene interaction
table. In practice that table is binary and tissue-blind. An edge asserts that a TF
\emph{can} bind a promoter, not that the TF is present in the cells generating the
measured signal. For a bulk time-series drawn from a whole organism this is a real
limitation, because the transcriptome is a mixture and the mixture changes.

Single-cell data is the obvious remedy, and iDREM already accepts it --- but only after
the fact. Its documentation is explicit that the cell-type data ``is not used when
predicting the iDREM model''; it annotates finished nodes rather than informing them.
The contribution of this work is to move that information upstream, into the model
itself, without modifying the engine. DREM reads the third column of its TF--gene file as
a score in $[0,1]$ rather than as a binary flag, and that column is a sufficient channel:
an edge whose TF is absent from the responding cell types can be down-weighted before the
search begins.

We test this on \textit{Arabidopsis thaliana} data from NASA's Open Science Data
Repository (OSDR). Screening all \NumArabidopsisStudiesScreened{} OSDR
\textit{Arabidopsis} studies for compatible genotype and treatment metadata yields two
cross-study pseudo-time-series. The first is a $\gamma$-irradiation series of
\NumCohortATimepoints{} timepoints spanning 10~min to 72~h. Its core --- OSD-498, OSD-508
and OSD-510 --- is a re-assembly of one experiment deposited under three accessions, that
of \citet{bourbousse2018sog1}, who built their own DREM model over these data and
reported \NumPublishedGroups{} co-expressed gene groups. That gives this work an unusual
advantage: a published DREM analysis of the same material to calibrate against, and a
genetic positive control, since the series contains both wild type and \textit{sog1-1},
a knockout of the NAC master regulator of the plant DNA-damage response
\citep{yoshiyama2009sog1}. The second series follows spaceflight seedling development
across \NumCohortBTimepoints{} timepoints and tests whether the method generalises to a
different stress.
